\naslov{Neenakomerne porazdelitve}

Težava pri neenakomernih porazdelitvah $Y$ je, da bomo za pogoste vrednosti
dobili manjšo napako.
Za porazdelitev
\[
  Y \sim
  \begin{pmatrix}
	0.01 & 0.99 \\
	\oplus & \ominus
  \end{pmatrix}
\]
je model $Y = \ominus$ perfekten, a ga ne želimo uporabiti iz očitnih razlogov.
Za diskreten $Y$ vzorčimo, torej upoštevamo manj velikih razredov
(\pojem{podvzorčenje}) ali pa primere majhnih razredov večkrat kopiramo
(\pojem{prevzorčenje}).

\vprasanje{Kaj je podvzorčenje in kaj prevzorčenje?}

Noben od teh pristopov ni idealen.
Malo boljši je algoritem SMOTE, kjer tvorimo nove sintetične primere.
Algoritem za izbran primer $(x,y)$ poišče $k$ najbližjih sosedov in za
naključnega $(x_n, y)$ z enakim $y$ tvori sintetični primer $(x + g \cdot (x_n -
x), y)$, kjer je $g \in [0,1]$ naključno izbran.
Dodatno algoritem sprejme tudi parametra o stopnji podvzorčenja in prevzorčenja;
to sta faktorja v številu novih razredov manjšinskega oz.~večinskega razreda
glede na originalno število.

\vprasanje{Opiši algoritem SMOTE.}

Pri regresiji neenakomerna porazdelitev izgleda tako, da imamo osamelce z zelo
različnimi vrednosti $Y$.
Definiramo \pojem{funkcijo pomembnosti} $\phi: Y \to [0,1]$, ki da visoko
pomembnost redkim primerom in zvezno pada k pogostejšim.
Primere z visoko pomembnostjo prevzorčimo, primere z nizko pomembnostjo pa
podvzorčimo.
V regresijski obliki algoritma SMOTE v sintetičnih primerih uporabimo
konveksno kombinacijo tudi v spremenljivki $Y$.

\vprasanje{Kako obravnavaš vprašanje o neenakomerni porazdelitvi v primeru
  regresije?}

% LocalWords:  podvzorčenje SMOTE podvzorčenja podvzorčimo
